% Français 9115, batch 19 — Gallica views 361–380.
% Pass: Opus 5.5 (claude-opus-5-5), 2026-09-22 — first pass, unchecked against the views by a human.
%
% What the twenty views hold. One piece, in French, in her hand: a fair copy
% of the latter part of her critique of Euler on the elastic lamina held by
% a stylet at an interior point — the same memoir whose working draft fills
% views 218–299 (batches 11–15), here recopied in a regular, well-spaced
% hand with almost no corrections, the paragraphs numbered 17 to 24 and
% worded, from § 18 on, as the draft's paragraphs of the same numbers
% (§ 19 at v. 283, § 20 at v. 287, § 21 at v. 291, § 22 at v. 297, § 23 at
% v. 299 of the draft). Where the draft had a crossed-out page (v. 289) or
% words between the lines, the fair copy has the text they settled on. It
% also departs from the draft in places, recorded in notes: at § 23 the
% one coefficient not zero for supported ends is written γ, and Euler's
% sign change is put on γ′ (γ = −γ′), where the draft was read δ and
% δ = −δ′; and Euler's values at x = 1/2 are cited to his § 48, where the
% draft cited § 50.
%
% Ten views carry text and ten are blank. Every leaf is written on the
% recto only and photographed on both sides: the written views are the odd
% ones, 361, 363 … 379, and the even ones, 362 … 380, are the blank backs,
% showing only the writing of the leaf before them reversed through the
% paper (checked by flopping crops of v. 366 and v. 372 and laying them on
% v. 365 and v. 371). There are no second exposures and no microfilm
% targets. The leaves are pasted on guards; the library's stamp is in the
% lower left margin of v. 361.
%
%   v. 361  § 17: the final equations for like ends are made of two terms,
%           each the product of the quantity that would vanish for a
%           separate lamina of length xa and of the one for (1−x)a; the
%           same analogy even for unlike ends (one fixed, one supported),
%           the stylet taken first as supported, then as fixed.
%   v. 363  end of § 17; § 18: yet the final equations yield little —
%           sin xω = 0 and its multiples for supported ends; for free or
%           fixed ends only the symmetry of the nodes, and x = 1/2.
%   v. 365  § 19: all of it follows without the final equations, from the
%           end conditions and y = 0 at the nodes; supported ends from
%           Euler's § 35; free ends from the equation of his § 21 (reading
%           uncertain), with x = 1/2 when sin ω is positive.
%   v. 367  the same for free ends, odd and even number of nodes, by
%           multiplying by 2 cos((1−2x)ω/2) and 2 sin((1−2x)ω/2) and
%           subtracting; that x and 1−x may be exchanged; the fixed-ends
%           equation beside the free-ends one.
%   v. 369  § 20: it was natural to examine Euler's method; his § 47,
%           where the stylet is supposed to break the continuity of the
%           curve and his equations I–IV are built on it; her objection
%           that the motion is one of the regular motions of the six cases,
%           hence continuous, hence α′ = α …; « Cependant je n'ose admettre
%           une opinion si contraire à celle d'Euler, sans la soumettre a
%           un examen plus approfondi ».
%   v. 371  § 21: the proposition — for like ends the two sets of
%           coefficients are equal — written in a larger, rounder hand;
%           that α = α′ entails the rest (equations III, IV); the case of
%           two free ends reduced to one identity.
%   v. 373  that identity developed term by term and verified for
%           sin ω positive (reference to her § 19).
%   v. 375  the same completed, then for sin ω negative.
%   v. 377  § 22, two fixed ends: the two expressions of α/α′ differ only
%           by the sign of e^{−xω}, and the change is compensated; § 23,
%           two supported ends: Euler's α/α′ = (e^{xω} − e^{(2−x)ω}) /
%           (e^{xω} − e^{−xω}) settles nothing because every coefficient
%           but one is zero; the eight coefficients he found for x = 1/2.
%   v. 379  Euler's reduction by sin ½ω; the conclusion that the stylet
%           gives nothing beyond y = 0 at its point combined with the end
%           conditions, the first two of his equations reducing to y = 0
%           and the last two to identities once α = α′, β = β′ are
%           admitted; § 24 opens (the half-laminae, odd number of nodes,
%           free and fixed ends) and stops mid-sentence on « dont une
%           extrémité ». This is past the last page of the draft read so
%           far (v. 299, § 23).
%
% Edges. View 361 opens a new paragraph (« 17. On voit que »); view 360
% is the blank back of the leaf before it (text reversed, a parenthesised
% page number at its head). The sentence at the foot of v. 379 continues on
% v. 381 (« et fixée et l'autre appuyée, des valeurs relatives au mouvement
% dela lame dont les … »), which carries « (31) » struck and « 186 »;
% v. 380 is blank. Views 360 and 381 were only glanced at, not transcribed.
%
% The numbers on the leaves, and why no \folio{} is written. Two series,
% one of each per written leaf, none on the blank backs. (1) At the head,
% in the middle, in her ink and in parentheses, each struck through with a
% horizontal stroke: (21) v. 361, (22) v. 363, (23) v. 365, (24) v. 367
% (struck twice, the 4 read with doubt), (25) v. 369, (26) v. 371, (27)
% v. 373, (28) v. 375, (29) v. 377, (30) v. 379 — her own pagination of
% the fair copy, cancelled. It is a different series from the draft's
% (underlined, not bracketed, and 30–38 at v. 283–299). (2) At the top
% right, in the broader pen and rounder hand of the running ink series:
% 176 (v. 361), 177, 178, 179, 180, 181, 182, 183, 184 (last figure
% repassed), 185 (v. 379, the 5 with a long descending tail), one to a
% leaf without a gap; v. 381 carries 186. It is penned, like the series
% the earlier batches recorded, so under this volume's house rule no
% \folio{} is written. Her paragraph numbers, written inside the first
% line: 17 (v. 361), 18 (v. 363), 19 (v. 365), 20 (v. 369), 21 (v. 371),
% 22 and 23 (v. 377), 24 (v. 379). Every number above was read on its
% view; none is inferred.
%
% Notation. Hers, kept: x a fraction of the length, the stylet at xa from
% one end and (1−x)a from the other; ω the angle of Euler's six cases;
% α β γ δ and α′ β′ γ′ δ′ his coefficients. Her γ is written like a « v »,
% her δ like a looped d. « sin », « cos », « tang » are written with a long
% s and are set \text{sin}, \text{cos}, \text{tang}; the stop she sometimes
% puts after them is kept where it stands. Braces { } are hers; where she
% writes two fractions side by side as one product (v. 371, 377) a \cdot is
% the transcription's. Roman numerals of Euler's equations are kept in
% capitals. Spelling is hers and not flagged: « extremités », « coëfficiens »,
% « stilet », « appuiée », « interromp », « inconcilliables »,
% « s'appercevoir », « problems », « mouvemens », « changemens », « cequé »,
% « cequi », « dela », « demême », « desorte », « parconséquent »,
% « aumilieu », « cidessus », « c'est adire », « demoitiés ». Where the leaf
% and the sense disagree the leaf stands, with a note: « sin xa = 0 »
% (v. 363), « (1−x)ω » for (1−x)a and « (1−x » unclosed (v. 365, 367).
%
% What could not be read. \ill{} stands in five places, all under a
% writer's heavy strike: three short words after « § 47 » at v. 369; a
% whole fraction after the first « = » at v. 371; one word after « § 47 »
% at v. 377; the first word of a struck group at v. 379. Readings offered
% with doubt are marked \uncertain{}: « § » before « precedent » (v. 361);
% the last figure of « § 21 » (v. 365); « §§ 15 » and « nœud » (v. 367);
% « résultats » (v. 367); « ordonnée qui », « dependantes » and the struck
% « les deux premieres » (v. 379). Nothing was guessed to fill a gap.
%
% Nothing here was checked against any published edition of Euler's memoir
% or of Sophie Germain's papers. No leaf of this batch bears on Fermat's
% Last Theorem; « Manuscript C » is not among these views.
\documentclass[11pt,a4paper]{article}
% The preamble shared by every transcription of the mathematicians' manuscripts in Gallica.
%
% It rests on one idea: **uncertainty compounds, it does not resolve.** A
% transcription that smooths over a doubtful word has destroyed the only thing
% it was made for — a reader must be able to tell, at every point, what was on
% the page and what was supplied. The apparatus macros below are therefore the
% critical apparatus, not decoration.
%
% The same commands are recognised by `scripts/render.mjs`, which produces the
% site's reading view, and by `scripts/tei.mjs`. A macro added here must be
% added there too, or rendering fails loudly — which is the wanted behaviour.
%
% Comments here are English, like the rest of the repository. The documents
% this preamble sets are French, and that is deliberate: see the skills.

\ProvidesPackage{germain}[2026/09/15 Transcriptions of mathematicians' manuscripts in Gallica]

\usepackage{amsmath,amssymb,amsthm}
\usepackage{mathrsfs}
% Kept from the parent preamble so the renderer's diagram support stays
% available; these manuscripts draw no commutative diagrams, and nothing here uses it.
\usepackage{tikz-cd}
\usepackage[english,french]{babel}
\usepackage{fontspec}

% --- Greek ------------------------------------------------------------------
% The text face stays fontspec's default, Latin Modern, which has no Greek:
% XeTeX drops every glyph it lacks with a line in the log and nothing on the
% page, and the Greek words the manuscripts quote vanished from the PDFs. So
% the Greek and Greek Extended blocks get a XeTeX character class of their own,
% and entering or leaving a run of it switches the family to CMU Serif — the
% same Computer Modern design, with polytonic Greek — and back. Only the
% family changes: a Greek word in \textit or \textbf keeps its shape and series.
% The transitions to and from every other class carry the tokens that class
% has towards an ordinary letter, so French spacing before « ; » or inside
% « » still holds after a Greek word. No grouping, so no brace or \endgroup
% can come out unbalanced; maths is left alone, since XeTeX inserts nothing
% there. Kept identical in `exercices.sty`.
\makeatletter
\newfontfamily\germainGreekfont{cmun}[
  Extension=.otf, NFSSFamily=germaingreek,
  UprightFont=*rm, ItalicFont=*ti, SlantedFont=*sl,
  BoldFont=*bx, BoldItalicFont=*bi, BoldSlantedFont=*bl]
\newXeTeXintercharclass\germain@greek
\def\germain@greekrange#1#2{%
  \@tempcnta=#1\relax
  \loop
    \XeTeXcharclass\@tempcnta=\germain@greek
    \ifnum\@tempcnta<#2\advance\@tempcnta\@ne\repeat}
\germain@greekrange{"0370}{"03FF}
\germain@greekrange{"1F00}{"1FFF}
\def\germain@greekprev{\rmdefault}
\def\germain@greekin{%
  \edef\germain@greekprev{\f@family}\fontfamily{germaingreek}\selectfont}
\def\germain@greekout{\fontfamily{\germain@greekprev}\selectfont}
\def\germain@greekjoin{%
  \XeTeXinterchartokenstate=\@ne
  \@tempcnta=\z@
  \loop
    \ifnum\@tempcnta=\germain@greek\else
      \XeTeXinterchartoks\germain@greek\@tempcnta=\expandafter{%
        \expandafter\germain@greekout\the\XeTeXinterchartoks\z@\@tempcnta}%
      \XeTeXinterchartoks\@tempcnta\germain@greek=\expandafter{%
        \the\XeTeXinterchartoks\@tempcnta\z@\germain@greekin}%
    \fi
    \ifnum\@tempcnta<4095\advance\@tempcnta\@ne\repeat}
% babel-french sets its own classes, and saves and restores the interchar
% state, each time a language is selected: join again after it does.
\addto\extrasfrench{\germain@greekjoin}
\addto\extrasenglish{\germain@greekjoin}
\AtBeginDocument{\germain@greekjoin}
\makeatother

\usepackage{geometry}
\geometry{a4paper,margin=25mm,marginparwidth=28mm}
\usepackage{marginnote}
\usepackage{xcolor}
\usepackage[normalem]{ulem}
\setlength{\emergencystretch}{3em}
\usepackage{hyperref}
\hypersetup{colorlinks=true,linkcolor=black,urlcolor=black,pdfborder={0 0 0}}

% --- Batch metadata ---------------------------------------------------------
% Taken as they stand by the rendering script, which shows them at the head of
% the reading view. They fix what the file claims to cover. The macro names are
% the parent project's, so that the scripts stay shared; \folder here means the
% volume's slug — `fr-9115` — and \pages counts Gallica views.

\newcommand{\folder}[1]{\gdef\grothFolder{#1}}
\newcommand{\batch}[1]{\gdef\grothBatch{#1}}
\newcommand{\pages}[2]{\gdef\grothFirst{#1}\gdef\grothLast{#2}}
\newcommand{\foldertitle}[1]{\gdef\grothFoldertitle{#1}}

% The holder's own shelfmark, printed as they write it: « Français 9115 ».
\newcommand{\shelfmark}[1]{\gdef\grothShelfmark{#1}}

% The dating, copied verbatim from the holder's catalogue — never inferred
% here. « XIXe siècle » is the BnF's; a pass that dates a leaf from its content
% says so in a \note{}, as its own inference.
\newcommand{\dating}[1]{\gdef\grothDating{#1}}

% The status notice, printed in the title block and repeated as a diagonal
% watermark on every page of the PDF. The manuscripts are in the public
% domain; what this line declares is not a right but a quality — an
% unverified first pass by a machine. The skills write the line; a file
% without it gets no watermark, so leaving it out is visible in the output.
\newcommand{\watermark}[1]{\gdef\grothWatermark{#1}}

\gdef\grothFolder{?}\gdef\grothBatch{}\gdef\grothFirst{?}\gdef\grothLast{?}%
\gdef\grothFoldertitle{}\gdef\grothShelfmark{}\gdef\grothDating{}\gdef\grothWatermark{}

% --- The résumé -------------------------------------------------------------
\newenvironment{resume}
  {\small\setlength{\parskip}{0.45em}}
  {\par\medskip\hrule\medskip}

% --- Keywords ---------------------------------------------------------------
% In English, closing the modernised reading's résumé: the modern vocabulary
% under which to search for what the leaves construct. The manifest extracts
% them to tag the volume — this line is the single source of the tags.
\newcommand{\keywords}[1]{\par\smallskip\noindent
  {\footnotesize\textsc{Keywords} --- \textit{#1}}\par}

% --- The critical apparatus -------------------------------------------------

% The Gallica view number, in the margin. This is the anchor that lets a
% disagreement over a formula be settled by looking at *one* image rather than
% a volume of 750. The site uses it to turn the facsimile as the transcription
% is scrolled. A view is one scanned image: usually one side of a leaf.
\newcommand{\page}[1]{%
  \marginnote{\footnotesize\color{gray!70}\textsf{v.\,#1}}%
  \label{page:#1}\ignorespaces}

% The BnF's pencilled foliation, where the leaf carries it and a pass can read
% it: « 198r ». This is what the literature cites — « ff. 198r–208v » — and
% the one line that turns a citation into a view. Recorded, never inferred:
% a folio a pass cannot read is not written.
\newcommand{\folio}[1]{%
  \marginnote{\footnotesize\color{gray!50}\textsf{f.\,#1}}\ignorespaces}

% The modernised reading's coarser counterpart to \page{}: which views a
% section is drawn from.
\newcommand{\pagerange}[2]{%
  \marginnote{\footnotesize\color{gray!70}\textsf{v.\,#1--#2}}%
  \ignorespaces}

% Illegible: never guessed.
\newcommand{\ill}{\textcolor{red!60!black}{[\,\ldots\,]}}

% A doubtful reading: the word is given, and flagged as uncertain.
\newcommand{\uncertain}[1]{\uline{#1}}

% An editorial addition — a missing letter, an implied word.
\newcommand{\add}[1]{\textcolor{blue!50!black}{[#1]}}

% Struck out by the writer. What was crossed out is part of the document.
\newcommand{\struck}[1]{\sout{#1}}

% The transcriber's note, on the state of the page or the mathematics.
\newcommand{\note}[1]{\par\smallskip{\small\color{gray!60!black}\textit{[#1]}}\par\smallskip}

% A marginal note of hers, distinct from the body of the text.
\newcommand{\marginal}[1]{\par\smallskip\noindent\hspace{1em}%
  {\small\color{gray!60!black}#1}\par\smallskip}

% --- Title ------------------------------------------------------------------
% The title block is French, like the documents themselves.

\AddToHook{shipout/background}{%
  \ifx\grothWatermark\empty\else
    \begin{tikzpicture}[remember picture, overlay]
      \node[rotate=52, scale=3.4, align=center, text=gray!18,
            font=\sffamily\bfseries]
        at (current page.center) {\grothWatermark};
    \end{tikzpicture}%
  \fi}

\AtBeginDocument{%
  \begin{center}
    {\large\bfseries Manuscrits de mathématiciens (Gallica) ---
      \ifx\grothShelfmark\empty\grothFolder\else\grothShelfmark\fi}\\[2pt]
    {\itshape\grothFoldertitle}\\[4pt]
    {\small vues \grothFirst--\grothLast
      \ifx\grothBatch\empty\else\ (lot \grothBatch)\fi}\\[2pt]
    \ifx\grothDating\empty\else
      {\small\color{gray!60!black}Datation du catalogue : \grothDating}\\[2pt]
    \fi
    \ifx\grothWatermark\empty\else
      {\small\bfseries\color{red!45!gray}\grothWatermark}\\[2pt]
    \fi
    {\footnotesize\color{gray!60!black}Transcription établie d'après les images IIIF de Gallica.
     Source gallica.bnf.fr / Bibliothèque nationale de France.\\
     Les lectures incertaines sont soulignées, les illisibles notées [\,\ldots\,],
     les ajouts entre crochets ; \textsf{v.} numérote les vues de Gallica,
     \textsf{f.} la foliotation au crayon de la BnF.}
  \end{center}
  \vspace{1em}}

\endinput


\folder{fr-9115}
\shelfmark{Français 9115}
\batch{19}
\pages{361}{380}
\dating{1801-1900}
\watermark{Lecture automatique\\non vérifiée}
\foldertitle{Papiers de Sophie GERMAIN. — Recueil de dissertations et problèmes mathématiques et physiques.}

\begin{document}

\note{Numérotation des feuillets. Chaque feuillet écrit porte deux nombres
à l'encre, et les versos blancs n'en portent aucun : au milieu de la tête,
entre parenthèses et barré d'un trait, sa pagination de cette mise au net,
(21) à la vue 361 puis (22) à (30) aux vues 363 à 379 ; en haut à droite,
d'une plume plus large et d'une écriture plus ronde, 176 à la vue 361 puis
177 à 185 aux vues 363 à 379, un par feuillet et sans lacune. Cette
seconde série occupe la place de la foliotation de la Bibliothèque, mais
elle est tracée à la plume, comme dans les lots précédents de ce volume ;
aucun « f. » n'est donc porté. Tous ces nombres ont été lus sur la vue ;
aucun n'est déduit. Les vues 362, 364, 366, 368, 370, 372, 374, 376, 378
et 380 sont les versos blancs des feuillets qui les précèdent et ne sont
pas transcrites.}

\page{361}

17. On voit que, dans le cas où l'état des extremités est semblable, les équations
finales sont composées de deux termes, dont l'un est le produit dela quantité
qui devroit être égalée a zéro si la partie dela lame de longueur $xa$ se
mouvoit comme une lame entière et separée dont une extremité seroit appuyée
par l'effet de l'application du stilet et dont l'autre resteroit libre fixée
ou appuiée suivant les cas; par la quantité qui devroit être égalée a zéro
si la partie de longueur $a(1-x)$ se mouvoit comme une lame entière et separée
qui auroit une de ses extremités fixée par l'effet de l'application du
stilet et l'autre libre fixée ou appuiée suivant les cas; et que le second
terme est composé d'une manière semblable puisqu'il ne differe du premier
que par le changement de $xa$ en $(1-x)a$ \emph{et vice versâ}.
\note{En tête, au milieu, « (21) » entre parenthèses, barré d'un trait
horizontal ; en haut à droite « 176 », à l'encre, d'une plume large. Le
paragraphe porte le numéro 17, écrit dans la ligne. « et vice versâ » est
souligné. C'est une mise au net : écriture régulière, presque sans
rature. Le feuillet est collé sur un onglet ; l'estampille de la
Bibliothèque impériale est apposée dans la marge de gauche, en bas.}

On peut encore observer que, même dans le cas où l'état des extremités est
different, il existe une grande analogie entre l'équation finale et celles
des cas précédens. en effet, si on choisit l'exemple d'une extremité fixée
et l'autre appuyée \uncertain{§} precedent, on voit que, si on suppose le stilet
simplement appuyé, l'une des quantités $\text{sin}\,x\omega$ . $\text{tang}\,(1-x)\omega - \frac{e^{(1-x)\omega} - e^{(x-1)\omega}}{e^{(1-x)\omega} + e^{(x-1)\omega}}$
est celle qui doit être égalée à zéro pour avoir le mouvement d'une lame
separée de longueur $xa$ dont une des extremités est appuiée par l'effet de
l'application du stilet et \add{dont} l'autre est également appuyée en vertu de l'hypothèse,
tandis que la seconde des mêmes quantités est celle qui doit être égalée
a zéro pour avoir le mouvement d'une lame separée de longueur $(1-x)a$
dont une des extremités est appuiée par l'effet de l'application du stilet
et dont l'autre est fixée en vertu de l'hypothèse; Et demême, si on
considère le stilet comme fixé, l'une des quantités $\text{tang}\,x\omega - \frac{e^{x\omega} - e^{-x\omega}}{e^{x\omega} + e^{-x\omega}}$
$\text{cos}\,(1-x)\omega - \frac{2}{e^{(1-x)\omega} + e^{(x-1)\omega}}$ sera celle qui doit être égalée à zéro pour avoir
la condition du mouvement d'une lame separée de longueur $xa$
\note{Le signe lu « § » devant « precedent » ressemble à un D pointé ;
la lecture est offerte d'après le sens. Après « tang », l'argument
$(1-x)\omega$ est écrit sur un autre tracé et empâté ; il se lit sous la
tache. Le point entre $\text{sin}\,x\omega$ et $\text{tang}\,(1-x)\omega$ sépare
les deux « quantités » dont parle la phrase (« l'une des quantités… la
seconde des mêmes quantités ») ; de même, plus bas, rien ne sépare sur le
feuillet $\text{tang}\,x\omega - \dots$ de $\text{cos}\,(1-x)\omega - \dots$, écrit
en tête de la ligne suivante. « dont » est ajouté au-dessus de la ligne.
La phrase se poursuit au feuillet suivant (vue 363).}

\page{363}

dont une des extremités seroit fixée en vertu de l'application du stilet et
dont l'autre seroit appuyée par l'hypothese relative aux extremités dela
lame primitive, tandis que la seconde des mêmes quantités seroit celle qui
devroit être egalée a zéro pour obtenir la condition du mouvement d'une
lame separée de longueur $(1-x)a$ dont une des extremités seroit fixée
par l'effet de l'application du stilet et l'autre pareillement fixée en
conséquence de l'hypothese.
\note{En tête, au milieu, « (22) » entre parenthèses, barré d'un trait
horizontal ; en haut à droite « 177 », à l'encre. Une petite tache
d'encre, en forme de C, au-dessus de « condition », n'est pas un mot.}

18. Cependant, malgré des analogies si remarquables, les équations
finales ne fournissent qu'un petit nombre de conséquences qui se réduisent
aux suivantes:

Pour le cas où les extremités sont l'une et l'autre appuyées, l'équation
finale montre que, si il y a un point d'appui aux distances $xa$ et
$(1-x)a$ des extremités, on aura $\text{sin}\,x\omega = 0$, et qu'il se formera encore
des points de repos à des distances $2xa$ et $(1-2x)a$, $3xa$ et $(1-3x)a$ \&c.
des mêmes extremités, parceque la condition $\text{sin}\,xa = 0$ entraine les
suivantes $\text{sin}\,2x\omega = 0$ $\text{sin}\,3x\omega = 0$ \&c; desorte que, pour une valeur
quelconque de $x$, on pourra toujours considérer la lame primitive
comme separée en deux autres lames qui auroient également leurs
deux extremités appuiées. Et, pour les cas des deux extremités libres
ou fixées, le seul resultat générale des équations finales est que,
si il existe un nœud de vibration à la distance $ax$ d'une des extremités,
il se formera un autre nœud à la même distance $ax$ de l'autre
extremité, parceque ces équations sont fonctions homogenes des
quantités $x$ et $x-1$. mais on déduit, à la verité, des mêmes équations
que les lames dont il s'agit peuvent etre considérées comme deux
lames separées, \struck{la} demoitiés moins longues, lorsque $x = \frac{1}{2}$.
\note{Le paragraphe porte le numéro 18, écrit dans la ligne. « $\text{sin}\,xa = 0$ »
est écrit ainsi, avec $a$ au lieu de $\omega$, deux lignes après
« $\text{sin}\,x\omega = 0$ » ; la leçon du feuillet est gardée. Le dernier
chiffre de « $x-1$ » est empâté. Le mot biffé devant « demoitiés » est
court et surchargé ; « la » est la lecture la plus probable, non sûre.
C'est la mise au net de la matière du § 18 du brouillon (voir la vue 283,
au lot 15, dont la première phrase reprend presque mot pour mot la fin
de l'alinéa sur les extrémités appuyées).}

\page{365}

19. Il me semble que l'on peut parvenir à tous ces resultats, sans avoir
recours aux équations finales, et en se bornant à la seule considération
des extremités et à la supposition $y = 0$ dans les points où la
courbe doit couper l'axe ou ligne de repos.
\note{En tête, au milieu, « (23) » entre parenthèses, barré d'un trait
horizontal ; en haut à droite « 178 », à l'encre. Le paragraphe porte le
numéro 19, écrit dans la ligne. La phrase d'ouverture est la mise au net
de celle qui ouvre le § 19 du brouillon, à la vue 283 (lot 15).}

En effet, pour le cas des deux extremités appuyées, on déduit des valeurs
données, Memoire d'Euler § 35, que l'ordonnée est en général exprimée par
l'équation $y = \dots\ \text{sin}\,x\omega$; desorte que, si par une cause quelconque la courbe
coupe l'axe dans le point qui répond à une certaine valeur de $x$, cequi
arrive pour les nœuds de vibration, on aura $\text{sin}\,x\omega = 0$, d'où on
déduira également $\text{sin}\,2x\omega = 0$ \&c., et parconséquent aussi que la lame
peut etre considerée comme autant de lames separées des longueurs
$xa$, ou comme deux lames egalement separées des longueurs $xa$ $(1-x)\omega$
dont les deux extremités seroient appuyées; car il est visible qu'à cause
de $\text{sin}\,\omega = 0$; $\text{sin}\,x\omega = 0$ entraine $\text{sin}\,(1-x)\omega = 0$.
\note{« $y = \dots\ \text{sin}\,x\omega$ » : les points sont sur le feuillet et
tiennent la place d'un facteur omis. « qui », après « le point », est
écrit sur un autre mot et empâté. « $(1-x)\omega$ », à la fin de
l'alinéa, est bien écrit avec $\omega$ là où le sens attend $(1-x)a$ ; la
leçon du feuillet est gardée. Dans « que la lame », l'article est
surchargé.}

A l'egard du cas où les extremités sont libres, lorsque la courbe
coupe l'axe, pour une valeur quelconque de $x$, on a l'équation
\[ e^{x\omega} \pm e^{(1-x)\omega} + (1 \mp e^{\omega})\,\text{sin}\,x\omega + (1 \pm e^{\omega})\,\text{cos}\,x\omega = 0 \]
(Voyez M. d'Euler § 2\uncertain{1});
et, en prenant d'abord les signes inferieurs qui conviennent au cas où
$\text{sin}\,\omega$ est positif, qui est celui où il y a un nombre impair de nœuds,
on trouve, en faisant $x = \frac{1}{2}$, $(1 + e^{\omega})\,\text{sin}\,\frac{1}{2}\omega + (1 - e^{\omega})\,\text{cos}\,\frac{1}{2}\omega = 0$ ou
$\text{tang}\,\frac{1}{2}\omega = \frac{e^{\omega} - 1}{e^{\omega} + 1}$, qui est la condition du mouvement d'une lame de
longueur $\frac{1}{2}a$, dont une extremité est appuyée et l'autre libre, et delaquelle
il resulte parconséquent que la lame dont il s'agit peut etre
considérée comme partagée en deux autres lames moitiée moins
longues, comme on l'a vu cidessus
\note{Le signe $\pm$ qui suit $e^{x\omega}$ est repassé et empâté. Le
dernier chiffre du renvoi « § 21 » est empâté ; il est offert avec doute.
« M. » est écrit ainsi, abrégé. Le feuillet s'arrête après « cidessus »,
sans ponctuation, sur le dernier tiers de la page laissé blanc.}

\page{367}

Si on reprend l'équation $(1 + e^{\omega})\,\text{sin}\,\frac{1}{2}\omega + (1 - e^{\omega})\,\text{cos}\,\frac{1}{2}\omega = 0$, qui resulte d'ailleurs des
conditions $\text{cos}.\,\omega = \frac{2e^{\omega}}{e^{2\omega} + 1}$ et $\text{sin}.\,\omega = \frac{e^{2\omega} - 1}{e^{2\omega} + 1}$ qui conviennent au cas present, \add{(Voy. M. d'Euler \uncertain{§§ 15})} parcequ'en effet,
lorsqu'il y a un nombre impair de nœuds, il y en a nécessairement un aumilieu
de la courbe, et si l'on multiplie cette équation par $2\,\text{cos}\,\frac{(1-2x)\omega}{2}$,
on aura
\[ 2\left\{(1 + e^{\omega})\,\text{sin}\,\tfrac{1}{2}\omega + (1 - e^{\omega})\,\text{cos}\,\tfrac{1}{2}\omega\right\}\text{cos}\,\frac{(1-2x)\omega}{2} \]
\[ = (1 + e^{\omega})\{\text{sin}.\,x\omega + \text{sin}\,(1-x)\omega\} + (1 - e^{\omega})\{\text{cos}\,x\omega + \text{cos}\,(1-x)\omega\} = 0 \]
et on trouvera, en retranchant l'équation generale de cette derniere,
\[ (1 + e^{\omega})\{\text{sin}\,x\omega + \text{sin}\,(1-x)\omega\} + (1 - e^{\omega})\{\text{cos}\,x\omega + \text{cos}\,(1-x)\omega\} \]
\[ - \{e^{x\omega} - e^{(1-x)\omega} + (1 + e^{\omega})\,\text{sin}\,x\omega + (1 - e^{\omega})\,\text{cos}\,x\omega\} \]
\[ = e^{(1-x)\omega} - e^{x\omega} + (1 + e^{\omega})\,\text{sin}\,(1-x)\omega + (1 - e^{\omega})\,\text{cos}\,(1-x)\omega = 0, \]
d'où on conclut que l'ordonnée est encore zéro lorsque l'on change $x$ et $1-x$,
\emph{et vice versâ}; et parconséquent que, si il y a un \uncertain{nœud} à la distance $ax$ d'une des
extremités, il y en aura necessairement un autre à la même distance de l'autre
extremité.
\note{En tête, au milieu, un nombre entre parenthèses, barré de deux
traits horizontaux ; sous les traits on lit « (24) », le second chiffre
avec doute ; en haut à droite « 179 », à l'encre. Le renvoi « (Voy. M.
d'Euler §§ 15) » est écrit au-dessus de la ligne, au-dessus de « present,
parcequ'en effet », et souligné d'un long trait ; les deux signes lus
« §§ » ressemblent à deux « s ». « et vice versâ » est souligné. Le mot
lu « nœud » est empâté par une tache.}

Demême aussi, pour le cas où il y a un nombre pair de nœuds, on
a l'équation $(e^{\omega} + 1)\,\text{sin}\,\frac{1}{2}\omega - (1 - e^{\omega})\,\text{cos}\,\frac{1}{2}\omega = 0$ qui resulte des conditions $\text{cos}.\,\omega = \frac{2e^{\omega}}{e^{2\omega} + 1}$
$\text{sin}.\,\omega = -\frac{e^{2\omega} - 1}{e^{2\omega} + 1}$; et on trouve, en la multipliant par $2\,\text{sin}\,\frac{(1-2x)\omega}{2}$,
\[ 2\left\{(e^{\omega} + 1)\,\text{sin}\,\tfrac{1}{2}\omega - (1 - e^{\omega})\,\text{cos}\,\tfrac{1}{2}\omega\right\}\text{sin}\,\frac{(1-2x)\omega}{2} \]
\[ = (e^{\omega} + 1)\{\text{cos}\,x\omega - \text{cos}\,(1-x)\omega\} + (1 - e^{\omega})\{\text{sin}.\,x\omega - \text{sin}\,(1-x)\omega\} = 0, \]
et cette équation retranchée de l'équation générale propre à ce cas, savoir de
\[ e^{x\omega} + e^{(1-x)\omega} + (1 - e^{\omega})\,\text{sin}.\,x\omega + (1 + e^{\omega})\,\text{cos}\,x\omega = 0, \]
donne
\[ e^{x\omega} + e^{(1-x)\omega} + (1 - e^{\omega})\,\text{sin}\,x\omega + (1 + e^{\omega})\,\text{cos}\,x\omega \]
\[ - \{(e^{\omega} + 1)[\text{cos}\,x\omega - \text{cos}\,(1-x)\omega] + (1 - e^{\omega})[\text{sin}.\,x\omega - \text{sin}\,(1-x)\omega]\} \]
\[ = e^{(1-x)\omega} + e^{x\omega} + (1 - e^{\omega})\,\text{sin}\,(1-x)\omega + (1 + e^{\omega})\,\text{cos}\,(1-x)\omega = 0; \]
d'où on tire, comme cidessus, la condition que $x$ et $(1-x$ peuvent etre changés
l'un en l'autre, sans que l'ordonnée cesse d'être egale a zéro
\note{Dans la seconde équation de l'alinéa, le signe entre
$\text{sin}.\,x\omega$ et $\text{sin}\,(1-x)\omega$ est une tache d'encre allongée ;
c'est un « $-$ » repassé, comme le montre la même accolade recopiée
deux lignes plus bas, où il est net. Dans l'équation générale et dans
la ligne qui la reprend, « $\text{cos}\,x\omega$ » est écrit sur un autre tracé
et empâté, aux deux endroits de la même manière. « $(1-x$ » est écrit
sans la parenthèse fermante.}

Les \uncertain{résultats} sont parfaitement semblables pour le cas où les deux extremités
sont fixés; car la seule différence entre ce cas et le precedent, est qu'on a ici
\[ e^{x\omega} \pm e^{(1-x)\omega} - \{(1 \mp e^{\omega})\,\text{sin}\,x\omega + (1 \pm e^{\omega})\,\text{cos}.\,x\omega\} = 0 \]
au lieu de
\[ e^{x\omega} \pm e^{(1-x)\omega} + (1 \mp e^{\omega})\,\text{sin}\,x\omega + (1 \pm e^{\omega})\,\text{cos}.\,x\omega\} \]
\note{« résultats » est surchargé : les premières lettres sont écrites sur
un autre mot. « au lieu de » est écrit dans la marge de gauche, en face
de la seconde formule. Celle-ci se ferme sur une accolade sans accolade
ouvrante, et le feuillet s'arrête là, sans ponctuation ; le feuillet
suivant (vue 369) ouvre un nouveau paragraphe, numéroté 20.}

\page{369}

20. Après avoir remarqué que les équations finales semblent indiquer des conclusions
qui ne sont pourtant pas admissibles puisqu'elles sont inconcilliables avec les
hypothèses subsistantes sur l'état des extremités, et après avoir prouvé que celles
des conséquences de ces équations qui sont applicables, peuvent être atteintes par
des considérations fort simples et independantes des mêmes équations, il étoit
naturel d'examiner la méthode au moyen delaquelle elles ont été trouvées.
\note{En tête, au milieu, « (25) » entre parenthèses, barré d'un trait
léger ; en haut à droite « 180 », à l'encre. Le paragraphe porte le
numéro 20, écrit dans la ligne. Dans « inconcilliables », les lettres du
milieu sont repassées. C'est la mise au net du § 20 du brouillon (vue
287, lot 15).}

On peut voir dans le § 47 \struck{\ill{}} du mémoire d'Euler, que l'auteur,
après avoir admis que le stilet soit simplement appuyé de manière à cequé
le mouvement puisse se communiquer de l'une à l'autre des deux parties
qu'il sépare, suppose cependant que l'action de ce stilet interromp la
continuité de la courbe, desorte que, suivant cette supposition, les deux
portions de cette courbe seroient exprimées par des équations qui differeroient
seulement par les valeurs des coëfficiens $\alpha\ \beta\ \gamma$ et $\delta$, et que c'est, d'après cette
hypothèse, qu'il a établi les équations I II III et IV
\note{Entre « § 47 » et « du mémoire », trois mots courts sont biffés
d'un trait épais qui les rend illisibles. Le $\gamma$ est tracé comme un
« v ». Les chiffres romains des équations sont tracés gras.}

Sans avoir recours au calcul, il peut d'abord paroître étonnant qu'Euler
suppose que les deux portions séparées par l'application du stilet, n'appartiennent
plus a une seule et même courbe; car le seul effet de cette application
est de déterminer la formation d'un certain nombre de nœuds, et parconséquent
de donner lieu à de certains mouvemens reguliers compris au nombre de
ceux qui ont été traités dans l'examen des six cas relatifs à l'état des
extremités; Or les mouvemens dont il s'agit s'exécutent suivant une courbe
continue: ces reflexions mènent donc à conclure que les coëfficiens $\alpha', \beta', \gamma', \delta'$
sont égaux chacun à son correspondant parmi les coëfficiens $\alpha, \beta, \gamma, \delta$.
Cependant je n'ose admettre une opinion si contraire à celle d'Euler, sans
la soumettre a un examen plus approfondi. c'est pourquoi je me propose
de démontrer la proposition suivante que je crois egalement facile à
\note{« portions » est souligné. « stilet », dans « l'application du
stilet », est écrit sur un autre mot. La phrase se poursuit au feuillet
suivant (vue 371). C'est la mise au net de la fin du brouillon de la
vue 289 (lot 15).}

\page{371}

démontrer pour les cas où l'état des extremités n'est pas le même pour chacune d'elles.
\note{En tête, au milieu, « (26) » entre parenthèses, barré d'un trait
horizontal ; en haut à droite « 181 », à l'encre. Cette première ligne
achève la phrase commencée au bas de la vue 369.}

21 Toutes les fois que la condition de l'application d'un stilet simplement
appuyé dans un point quelconque de l'étendue d'une lame elastique vibrante,
laisse subsister une hypothèse déterminée et semblable sur l'état de l'une
et l'autre des extremités de la même lame, les coëfficiens $\alpha, \beta, \gamma\ \delta$ et
$\alpha', \beta', \gamma'\ \delta'$ qui entrent dans l'expression des ordonnées des deux parties
de la courbe séparées par ce point d'appui, sont égaux chacun à son
correspondant.
\note{Le paragraphe porte le numéro 21. L'énoncé de la proposition, de
« Toutes les fois » à « correspondant », est écrit d'une écriture plus
droite, plus ronde et plus détachée que le reste du feuillet, comme pour
le mettre en valeur ; la suite (« Je vais commencer ») revient à la main
courante. Rien sur le feuillet ne dit s'il s'agit d'une autre main ; la
passe n'en décide pas.}

Je vais commencer par prouver qu'en général il suffit que l'un de ces coëfficiens
soit égal à son correspondant, pour que la même égalité ait lieu entre tous les
autres, desorte que, si on a, par exemple, $\alpha = \alpha'$, il en resultera nécessairement
$\beta = \beta'$ $\gamma = \gamma'$ \&c.

Je suppose donc $\alpha = \alpha'$; il resulte alors de l'équation $\alpha e^{x\omega} + \beta e^{-x\omega} = \alpha' e^{x\omega} + \beta' e^{-x\omega}$
§. 8 que $\beta = \beta'$. dans cette hypothèse, les équations III et IV se reduisent à
\[ \text{III}.\ (\gamma - \gamma')\,\text{cos}\,x\omega - (\delta - \delta')\,\text{sin}.\,x\omega = 0 \qquad \text{IV}.\ (\gamma - \gamma')\,\text{sin}.\,x\omega + (\delta - \delta')\,\text{cos}.\,x\omega = 0 \]
on tire de la première multipliée par $\text{cos}\,x\omega$ et ajoutée à la seconde multipliée elle même
par $\text{sin}\,x\omega$, $(\gamma - \gamma')(\text{sin}^2 x\omega + \text{cos}^2 x\omega) = 0$ c'est adire $\gamma = \gamma'$; on a alors
\[ \text{III}.\ (\delta - \delta')\,\text{sin}\,x\omega = 0 \qquad \text{IV}.\ (\delta - \delta')\,\text{cos}.\,x\omega = 0, \]
et parconséquent $\delta = \delta'$.
\note{« §. 8 », en tête de ligne, renvoie à l'équation qui précède ; le
signe du paragraphe est tracé comme un I bouclé. Les mots « par
$\text{sin}\,x\omega$ » débordent dans la marge de gauche. Le $\gamma$ est
tracé comme un « v » et les chiffres romains sont tracés gras, comme au
feuillet précédent.}

Ainsi il ne s'agit plus que de prouver qu'en effet la condition $\alpha = \alpha'$ est
nécessaire. je choisis d'abord le cas où les deux extremités sont supposées libres;
on a alors
\[ \frac{\alpha e^{-\omega}\{e^{(1-x)\omega} + \text{cos}\,(1-x)\omega + \text{sin}\,(1-x)\omega\}}{\alpha'\{e^{-x\omega} - \text{sin}\,x\omega + \text{cos}\,x\omega\}} = \]
\struck{\ill{}} (Voyez § 8)
\[ = \frac{(e^{(1-x)\omega} + e^{(x-1)\omega})\,\text{sin}\,(1-x)\omega - (e^{(1-x)\omega} - e^{(x-1)\omega})\,\text{cos}\,(1-x)\omega}{(e^{x\omega} - e^{-x\omega})\,\text{cos}\,x\omega - (e^{x\omega} + e^{-x\omega})\,\text{sin}.\,x\omega} \]
il suffit donc de montrer que l'équation
\[ \{e^{-x\omega} + e^{-\omega}(\text{cos}\,(1-x)\omega + \text{sin}\,(1-x)\omega)\}\{(e^{x\omega} + e^{-x\omega})\,\text{sin}\,x\omega - (e^{x\omega} - e^{-x\omega})\,\text{cos}\,x\omega\} \]
\[ = \{\text{sin}\,x\omega - \text{cos}.\,x\omega - e^{-x\omega}\}\{(e^{(1-x)\omega} + e^{(x-1)\omega})\,\text{sin}\,(1-x)\omega - (e^{(1-x)\omega} - e^{(x-1)\omega})\,\text{cos}\,(1-x)\omega\} \]
\note{Après le premier signe $=$, une fraction entière est biffée lettre
par lettre et n'est pas lue ; on devine seulement qu'elle commence au
numérateur et au dénominateur par un « 2 » et qu'elle contient
$e^{(1-x)\omega}$ en haut et $e^{x\omega}$, $e^{-x\omega}$ en bas. Le renvoi
« (Voyez § 8) » est écrit à sa droite, non biffé. La seconde fraction est
écrite en dessous, précédée de $=$. Le feuillet s'arrête sur la seconde
ligne de l'équation ; la phrase se poursuit au feuillet suivant (vue
373).}

\page{373}

a lieu d'elle même sous les conditions indiquées.
\note{En tête, au milieu, « (27) » entre parenthèses, barré d'un trait
horizontal ; en haut à droite « 182 », à l'encre. La première ligne
achève la phrase du bas de la vue 371 (« il suffit donc de montrer que
l'équation … a lieu d'elle même »). Le feuillet est presque tout entier
occupé par le calcul, disposé ligne à ligne ; la disposition des lignes
est gardée.}

On a par le développement
\[ (e^{(x-1)\omega} + e^{-(x+1)\omega})\{\text{cos}\,(1-x)\omega\,\text{sin}\,x\omega + \text{sin}\,(1-x)\omega\,\text{sin}\,x\omega\} \]
\[ - (e^{(x-1)\omega} - e^{-(x+1)\omega})\{\text{cos}\,(1-x)\omega\,\text{cos}\,x\omega + \text{sin}\,(1-x)\omega\,\text{cos}\,x\omega\} \]
\[ + e^{-x\omega}\{(e^{(1-x)\omega} + e^{(x-1)\omega})\,\text{sin}\,(1-x)\omega - (e^{(1-x)\omega} - e^{(x-1)\omega})\,\text{cos}\,(1-x)\omega\} \]
\[ = (e^{(1-x)\omega} + e^{(x-1)\omega})\{\text{sin}\,(1-x)\omega\,\text{sin}.\,x\omega - \text{sin}\,(1-x)\omega\,\text{cos}.\,x\omega\} \]
\[ - (e^{(1-x)\omega} - e^{(x-1)\omega})\{\text{cos}\,(1-x)\omega\,\text{sin}\,x\omega - \text{cos}\,(1-x)\omega\,\text{cos}\,x\omega\} \]
\[ - e^{-x\omega}\{(e^{x\omega} + e^{-x\omega})\,\text{sin}.\,x\omega - (e^{x\omega} - e^{-x\omega})\,\text{cos}.\,x\omega\} \]
ou
\[ e^{(x-1)\omega}\{\text{cos}\,(1-x)\omega\,\text{sin}\,x\omega + \text{sin}\,(1-x)\omega\,\text{sin}.\,x\omega - \text{cos}\,(1-x)\omega\,\text{cos}.\,x\omega - \text{sin}\,(1-x)\omega\,\text{cos}.\,x\omega\} \]
\[ + e^{-(1+x)\omega}\{\text{cos}\,(1-x)\omega\,\text{sin}.\,x\omega + \text{sin}\,(1-x)\omega\,\text{sin}.\,x\omega + \text{cos}\,(1-x)\omega\,\text{cos}.\,x\omega + \text{sin}\,(1-x)\omega\,\text{cos}\,x\omega\} \]
\[ + e^{-x\omega}\{(e^{(1-x)\omega} + e^{(x-1)\omega})\,\text{sin}\,(1-x)\omega - (e^{(1-x)\omega} - e^{(x-1)\omega})\,\text{cos}\,(1-x)\omega \]
\[ + (e^{x\omega} + e^{-x\omega})\,\text{sin}.\,x\omega - (e^{x\omega} - e^{-x\omega})\,\text{cos}.\,x\omega\} \]
\[ = e^{(x-1)\omega}\{\text{sin}.\,(1-x)\omega\,\text{sin}.\,x\omega - \text{sin}\,(1-x)\omega\,\text{cos}.\,x\omega + \text{cos}\,(1-x)\omega\,\text{sin}\,x\omega - \text{cos}\,(1-x)\omega\,\text{cos}\,x\omega\} \]
\[ + e^{(1-x)\omega}\{\text{sin}\,(1-x)\omega\,\text{sin}.\,x\omega - \text{sin}\,(1-x)\omega\,\text{cos}\,x\omega - \text{cos}\,(1-x)\omega\,\text{sin}.\,x\omega + \text{cos}\,(1-x)\omega\,\text{cos}\,x\omega\} \]
en effaçant ce qui se détruit, et multipliant par $e^{x\omega}$ cette équation devient
\[ e^{-\omega}\{\text{cos}\,(1-x)\omega\,\text{sin}.\,x\omega + \text{sin}.\,(1-x)\omega\,\text{sin}.\,x\omega + \text{cos}\,(1-x)\omega\,\text{cos}.\,x\omega + \text{sin}\,(1-x)\omega\,\text{cos}\,x\omega\} \]
\[ + (e^{(1-x)\omega} + e^{(x-1)\omega})\,\text{sin}\,(1-x)\omega - (e^{(1-x)\omega} - e^{(x-1)\omega})\,\text{cos}\,(1-x)\omega \]
\[ + (e^{x\omega} + e^{-x\omega})\,\text{sin}\,x\omega - (e^{x\omega} - e^{-x\omega})\,\text{cos}\,x\omega \]
\[ = e^{\omega}\{\text{sin}\,(1-x)\omega\,\text{sin}.\,x\omega - \text{sin}\,(1-x)\omega\,\text{cos}\,x\omega - \text{cos}\,(1-x)\omega\,\text{sin}\,x\omega + \text{cos}\,(1-x)\omega\,\text{cos}.\,x\omega\} \]
qui peut etre mise sous la forme
\[ e^{-\omega}\{\text{cos}\,(1-x)\omega + \text{sin}\,(1-x)\omega\}\{\text{cos}.\,x\omega + \text{sin}\,x\omega\} \]
\[ + e^{(1-x)\omega}\{\text{sin}\,(1-x)\omega - \text{cos}\,(1-x)\omega\} + e^{(x-1)\omega}\{\text{sin}\,(1-x)\omega + \text{cos}\,(1-x)\omega\} \]
\[ + e^{x\omega}\{\text{sin}\,x\omega - \text{cos}\,x\omega\} + e^{-x\omega}\{\text{sin}.\,x\omega + \text{cos}\,x\omega\} \]
\[ = e^{\omega}\{\text{sin}\,(1-x)\omega - \text{cos}\,(1-x)\omega\}\{\text{sin}\,x\omega - \text{cos}.\,x\omega\} \]
je suppose d'abord $\text{sin}.\,\omega$ positif, on a pour ce cas
\[ \left.\begin{array}{l} e^{x\omega} - e^{(1-x)\omega} + (1 + e^{\omega})\,\text{sin}\,x\omega + (1 - e^{\omega})\,\text{cos}.\,x\omega = 0 \\ e^{(1-x)\omega} - e^{x\omega} + (1 + e^{\omega})\,\text{sin}\,(1-x)\omega + (1 - e^{\omega})\,\text{cos}\,(1-x)\omega = 0 \end{array}\right\} \]
(voyez § 19)

et on tire du produit de ces équations multipliée par $e^{-\omega}$
\[ e^{-\omega}\{\text{cos}\,(1-x)\omega + \text{sin}\,(1-x)\omega\}\{\text{cos}\,x\omega + \text{sin}\,x\omega\} = e^{\omega}\{\text{sin}\,(1-x)\omega - \text{cos}\,(1-x)\omega\}\{\text{sin}\,x\omega - \text{cos}\,x\omega\} \]
\[ + 2 - e^{(1-2x)\omega} - e^{(2x-1)\omega} + e^{x\omega}\{\text{sin}\,(1-x)\omega - \text{cos}\,(1-x)\omega\} \]
\[ - e^{(1-x)\omega}\{\text{sin}\,(1-x)\omega - \text{cos}\,(1-x)\omega\} + e^{(1-x)\omega}\{\text{sin}\,x\omega - \text{cos}\,x\omega\} - e^{x\omega}\{\text{sin}\,x\omega - \text{cos}\,x\omega\} \]
\note{Le renvoi « (voyez § 19) » est écrit à droite de l'accolade qui
réunit les deux équations. Au bas du feuillet, le dernier signe $-$,
devant $e^{x\omega}$, est empâté. Le calcul se poursuit au feuillet
suivant (vue 375).}

\page{375}

on trouve ainsi, en retranchant cette équation de l'équation precedente,
\[ 2 - e^{(1-2x)\omega} - e^{(2x-1)\omega} + e^{x\omega}\{\text{sin}\,(1-x)\omega - \text{cos}\,(1-x)\omega\} \]
\[ - e^{(1-x)\omega}\{\text{sin}\,(1-x)\omega - \text{cos}\,(1-x)\omega\} + e^{(1-x)\omega}\{\text{sin}\,x\omega - \text{cos}\,x\omega\} - e^{x\omega}\{\text{sin}\,x\omega - \text{cos}\,x\omega\} \]
\[ + e^{(x-1)\omega}\{\text{sin}\,(1-x)\omega + \text{cos}\,(1-x)\omega\} + e^{-x\omega}\{\text{sin}\,x\omega + \text{cos}\,x\omega\} \]
\[ + e^{(1-x)\omega}\{\text{sin}\,(1-x)\omega - \text{cos}\,(1-x)\omega\} + e^{x\omega}\{\text{sin}\,x\omega - \text{cos}\,x\omega\} \]
\[ = 2 - e^{(1-2x)\omega} - e^{(2x-1)\omega} + e^{x\omega}\{\text{sin}\,(1-x)\omega - \text{cos}\,(1-x)\omega\} + e^{(x-1)\omega}\{\text{sin}\,(1-x)\omega + \text{cos}\,(1-x)\omega\} \]
\[ + e^{(1-x)\omega}\{\text{sin}\,x\omega - \text{cos}\,x\omega\} + e^{-x\omega}\{\text{sin}\,x\omega + \text{cos}\,x\omega\} \]
\[ = e^{-x\omega}\{e^{x\omega} - e^{(1-x)\omega} + (1 + e^{\omega})\,\text{sin}\,x\omega + (1 - e^{\omega})\,\text{cos}\,x\omega\} \]
\[ + e^{(x-1)\omega}\{e^{(1-x)\omega} - e^{x\omega} + (1 + e^{\omega})\,\text{sin}\,(1-x)\omega + (1 - e^{\omega})\,\text{cos}\,(1-x)\omega\} \]
quantité qui est nécessairement égale à zéro, puisqu'elle est la somme des
deux valeurs d'$y = 0$ qui répondent a $x$ et $x-1$, multipliées l'une par $e^{-x\omega}$,
et l'autre par $e^{(x-1)\omega}$.
\note{En tête, au milieu, « (28) » entre parenthèses, barré d'un trait
horizontal ; en haut à droite « 183 », à l'encre. La première ligne
enchaîne sur le calcul du bas de la vue 373 (« cette équation » est
celle qu'on vient d'y tirer du produit des deux équations). Le « 2 »
qui ouvre la troisième ligne de calcul est tracé gras, comme un $x$.}

Il faut maintenant prouver la même chose pour le cas où $\text{sin}.\,\omega$ est négatif.
les deux équations relatives a ce cas sont
\[ \left.\begin{array}{l} e^{x\omega} + e^{(1-x)\omega} + (1 - e^{\omega})\,\text{sin}.\,x\omega + (1 + e^{\omega})\,\text{cos}\,x\omega = 0 \\ e^{x\omega} + e^{(1-x)\omega} + (1 - e^{\omega})\,\text{sin}\,(1-x)\omega + (1 + e^{\omega})\,\text{cos}\,(1-x)\omega = 0 \end{array}\right\} \]
§ 19

et on tire de leur produit multiplié par $e^{-\omega}$
\[ e^{-\omega}\{\text{sin}\,(1-x)\omega + \text{cos}\,(1-x)\omega\}\{\text{sin}\,x\omega + \text{cos}\,x\omega\} = e^{\omega}\{\text{sin}\,(1-x)\omega - \text{cos}\,(1-x)\omega\}\{\text{sin}\,x\omega - \text{cos}\,x\omega\} \]
\[ + 2 + e^{(2x-1)\omega} + e^{(1-2x)\omega} \]
\[ - (e^{x\omega} + e^{(1-x)\omega})\{\text{sin}\,(1-x)\omega - \text{cos}\,(1-x)\omega + \text{sin}\,x\omega - \text{cos}\,x\omega\} \]
puis en retranchant cette équation de celle qu'il faut démontrer, on trouve
\[ 2 + e^{(2x-1)\omega} + e^{(1-2x)\omega} - (e^{x\omega} + e^{(1-x)\omega})\{\text{sin}\,(1-x)\omega - \text{cos}\,(1-x)\omega + \text{sin}\,x\omega - \text{cos}\,x\omega\} \]
\[ + e^{(x-1)\omega}\{\text{sin}\,(1-x)\omega + \text{cos}\,(1-x)\omega\} + e^{-x\omega}\{\text{sin}\,x\omega + \text{cos}\,x\omega\} \]
\[ + e^{(1-x)\omega}\{\text{sin}\,(1-x)\omega - \text{cos}\,(1-x)\omega\} + e^{x\omega}\{\text{sin}\,x\omega - \text{cos}\,x\omega\} \]
\[ = 2 + e^{(2x-1)\omega} + e^{(1-2x)\omega} - e^{x\omega}\{\text{sin}\,(1-x)\omega - \text{cos}\,(1-x)\omega\} + e^{(x-1)\omega}\{\text{sin}\,(1-x)\omega + \text{cos}\,(1-x)\omega\} \]
\[ - e^{(1-x)\omega}\{\text{sin}\,x\omega - \text{cos}\,x\omega\} + e^{-x\omega}\{\text{sin}\,x\omega + \text{cos}\,x\omega\} \]
\[ = e^{-x\omega}\{e^{x\omega} + e^{(1-x)\omega} + (1 - e^{\omega})\,\text{sin}\,x\omega + (1 + e^{\omega})\,\text{cos}\,x\omega\} \]
\[ + e^{(x-1)\omega}\{e^{x\omega} + e^{(1-x)\omega} + (1 - e^{\omega})\,\text{sin}\,(1-x)\omega + (1 + e^{\omega})\,\text{cos}\,(1-x)\omega\} \]
quantité qui par les raisons déduites cidessus, est nécessairement égale a zéro.
\note{Le renvoi « § 19 » est écrit à droite de l'accolade, sans
parenthèses. Les deux équations du cas où $\text{sin}\,\omega$ est négatif
commencent l'une et l'autre par « $e^{x\omega} + e^{(1-x)\omega}$ », qui
est symétrique. Dans la première ligne du produit, le signe à
l'intérieur de la première accolade est petit et empâté ; il se lit
« $+$ ». Le feuillet s'arrête sur « égale a zéro. » ; le quart inférieur
de la page reste blanc.}

\page{377}

22. Lorsque les extremités sont l'une et l'autre fixées, il est également
vrai que les deux series des coëfficiens $\alpha, \beta, \gamma, \delta$ et $\alpha', \beta', \gamma', \delta'$ sont égales entr'elles.
en effet, je me suis assurée que l'on a, pour cedernier cas,
\[ \frac{\alpha}{\alpha'} = \frac{\{(e^{(1-x)\omega} + e^{(x-1)\omega})\,\text{sin}\,(1-x)\omega - (e^{(1-x)\omega} - e^{(x-1)\omega})\,\text{cos}\,(1-x)\omega\}}{(e^{x\omega} + e^{-x\omega})\,\text{sin}\,x\omega - (e^{x\omega} - e^{-x\omega})\,\text{cos}\,x\omega} \]
\[ \cdot \frac{\{\text{sin}\,x\omega - \text{cos}.\,x\omega + e^{-x\omega}\}}{e^{-\omega}\{\text{sin}\,(1-x)\omega + \text{cos}\,(1-x)\omega\} - e^{-x\omega}} \]
au lieu de
\[ \frac{\alpha}{\alpha'} = \frac{\{(e^{(1-x)\omega} + e^{(x-1)\omega})\,\text{sin}\,(1-x)\omega - (e^{(1-x)\omega} - e^{(x-1)\omega})\,\text{cos}\,(1-x)\omega\}}{(e^{x\omega} + e^{-x\omega})\,\text{sin}\,x\omega - (e^{x\omega} - e^{-x\omega})\,\text{cos}\,x\omega} \]
\[ \cdot \frac{\{\text{sin}\,x\omega - \text{cos}\,x\omega - e^{-x\omega}\}}{e^{-\omega}\{\text{sin}\,(1-x)\omega + \text{cos}\,(1-x)\omega\} + e^{-x\omega}} \]
que j'ai montré avoir lieu pour le cas des extremités libres. il est aisé de
voir que la seule différence qui existe entre ces deux équations, consiste
dans le changement de signe de la quantité $e^{-x\omega}$; et, comme on a,
pour le cas des extremités fixées, $e^{x\omega} \mp e^{(1-x)\omega} - \{(1 \pm e^{\omega})\,\text{sin}\,x\omega + (1 \mp e^{\omega})\,\text{cos}\,x\omega\} = 0$,
au lieu de \dots\ $e^{x\omega} \mp e^{(1-x)\omega} + (1 \pm e^{\omega})\,\text{sin}\,x\omega + (1 \mp e^{\omega})\,\text{cos}\,x\omega = 0$
pour celui des extremités libres, il y a compensation entre les changemens
de signes; et on parvient ici, comme dans le cas precedent, à conclure
$\alpha = \alpha'$.
\note{En tête, au milieu, « (29) » entre parenthèses, barré d'un trait
horizontal ; en haut à droite « 184 », à l'encre, le dernier chiffre
repassé. Le paragraphe porte le numéro 22. Chaque valeur de
$\frac{\alpha}{\alpha'}$ est écrite, comme à la vue 371, en deux
fractions côte à côte sans signe entre elles ; le point de
multiplication est de la transcription. « au lieu de » est suivi, sur
le feuillet, d'une ligne de points qui aligne la seconde équation sous
la première. C'est la mise au net du § 22 du brouillon (vue 297,
lot 15).}

23. A l'égard du cas où l'une et l'autre extremités sont appuyées, Euler
trouve (Voyez § 47 \struck{\ill{}} de son memoire) $\frac{\alpha}{\alpha'} = \frac{e^{x\omega} - e^{(2-x)\omega}}{e^{x\omega} - e^{-x\omega}}$; et il est clair que cette équation
ne mène pas à conclure $\alpha = \alpha'$: mais il faut remarquer qu'ici tous les coëfficiens,
a l'exception de $\underline{\gamma}$, sont nuls en vertu de l'état des extremités; desorte que
$\frac{\alpha}{\alpha'}$ est essentiellement une quantité indéterminée; et que parconséquent on ne
peut rien conclure de l'équation $\frac{\alpha}{\alpha'} = \frac{e^{x\omega} - e^{(2-x)\omega}}{e^{x\omega} - e^{-x\omega}}$: c'est sans doute cette circonstance
qui a empêché l'auteur de s'appercevoir de l'identité entre les deux series
de coëfficiens. Cependant, après avoir trouvé, dans la supposition $x = \frac{1}{2}$ § 48,
\[ \alpha = e^{\frac{1}{2}\omega} - e^{\frac{3}{2}\omega};\quad \beta = e^{\frac{3}{2}\omega} - e^{\frac{1}{2}\omega}\quad \gamma = \frac{(e^{\frac{3}{2}\omega} - e^{\frac{1}{2}\omega})(e^{\frac{1}{2}\omega} - e^{-\frac{1}{2}\omega})}{\text{sin}\,\frac{1}{2}\omega},\quad \delta = 0 \]
\[ \alpha' = e^{\frac{1}{2}\omega} - e^{-\frac{1}{2}\omega}\quad \beta' = -e^{2\omega}(e^{\frac{1}{2}\omega} - e^{-\frac{1}{2}\omega}),\quad \gamma' = \frac{(e^{\frac{3}{2}\omega} - e^{\frac{1}{2}\omega})(e^{\frac{1}{2}\omega} - e^{-\frac{1}{2}\omega})}{\text{sin}\,\frac{1}{2}\omega - \text{cos}\,\frac{1}{2}\omega\,\text{tang}.\,\omega} \]
\[ \quad \delta' = \frac{(e^{\frac{1}{2}\omega} - e^{-\frac{1}{2}\omega})(e^{\frac{1}{2}\omega} - e^{\frac{3}{2}\omega})\,\text{tang}.\,\omega}{\text{sin}\,\frac{1}{2}\omega - \text{cos}\,\frac{1}{2}\omega\,\text{tang}.\,\omega} \]
\note{Le paragraphe porte le numéro 23. Après « § 47 », un mot court est
biffé d'un trait épais et n'est pas lu. La lettre soulignée après
« a l'exception de » est tracée comme son $\gamma$, un « v » bouclé ; le
brouillon correspondant (vue 299, lot 15) avait été lu $\delta$ au même
endroit, et le relevé des coefficients qui suit, où $\delta = 0$, s'accorde
avec $\gamma$. Dans « § 48 », le 4 est écrit sur un autre chiffre. Les
exposants lus $\frac{3}{2}$ ont un 3 à longue queue qui ressemble à un 9,
comme dans la copie d'Euler. Au début de la barre de fraction de $\gamma$,
un gros empâtement allongé ; il n'est pas lu comme un signe $-$, le
numérateur commençant ici par $e^{\frac{3}{2}\omega}$. Le $\gamma$ est
partout tracé comme un « v ». Le feuillet s'arrête sur ces huit valeurs ;
la phrase (« après avoir trouvé… ») se poursuit au feuillet suivant (vue
379). C'est la mise au net du § 23 du brouillon (vue 299).}

\page{379}

il multiplie tous ces coëfficiens par $\text{sin}\,\frac{1}{2}\omega$; cequi lui donne
\[ \alpha = \alpha' = 0\quad \beta = \beta' = 0\quad \gamma = -\gamma'.\ \text{et}\ \delta = \delta' = 0 \]
conformement à cequé j'ai dit devoir
avoir lieu toutes les fois que l'on a égard à l'état des extremités. on peut
observer seulement que le changement de signe est attribué au coëfficient $\underline{\gamma'}$, au lieu
d'être considéré comme appartenant à l'ordonnée dans l'expression delaquelle il
doit entrer, \uncertain{ordonnée qui} est prise de l'autre côté de l'axe.
\note{En tête, au milieu, « (30) » entre parenthèses, barré d'un trait
horizontal ; en haut à droite « 185 », à l'encre, le 5 à longue queue
descendante. La première ligne achève la phrase du bas de la vue 377
(« après avoir trouvé … »). « cequé » est écrit ainsi. Les mots lus
« ordonnée qui » sont écrits sur d'autres lettres et empâtés ; la
lecture est offerte d'après le sens. Une tache d'encre ronde couvre la
fin de « l'axe ». Le brouillon (vue 299, lot 15) avait ici
$\delta = -\delta'$ ; la mise au net écrit $\gamma = -\gamma'$ et nomme
$\gamma'$, souligné, comme le coefficient auquel Euler attribue le
changement de signe.}

On voit donc que l'examen de l'effet de l'application d'un stilet
dans un point quelconque de la longueur d'une lame elastique vibrante,
ne donne rien de plus que la supposition $y = 0$ dans cemême point
combiné avec les conditions \uncertain{dependantes} de l'état des extremités. aussi les
deux premières des équations I II III IV, employées à la solution des
problems de ce genre, se reduisent-elles \struck{\ill{} \uncertain{les deux premieres}} à
l'expression d'$y$ egalée à zéro, et les deux dernières à des équations identiques,
lorsque l'on admet $\alpha = \alpha'$ $\beta = \beta'$ \&c.
\note{« dependantes » est surchargé : les lettres du milieu sont écrites
sur d'autres. Le groupe biffé devant « à l'expression » est barré lettre
par lettre ; on croit y lire, après un premier mot illisible, « les deux
premieres ». « problems » est écrit ainsi. Les chiffres romains des
équations sont tracés gras.}

24. J'ai déjà remarqué que, dans le cas où il y a un nombre impair
de nœuds, la lame dont les deux extrémités sont libres, peut être considérée
comme partagée en deux lames séparées moitié moins longues qui auroient
une de leurs extrémités libre et l'autre appuiée; et que demême aussi la lame
dont les deux extrémités sont fixées peut être considerée dans lemême cas
comme partagée en deux lames séparées moitié moins longues qui auroient
une de leurs extrémités fixée et l'autre appuiée. Cette remarque autorise
à conclure toutes les circonstances du mouvement d'une lame dont une
extrémité est libre et l'autre appuiée, de celles qui appartiennent au
mouvement dela lame dont les deux extrémités sont libres, et demême
encore cequi concerne le mouvement dela lame dont une extrémité
\note{Le paragraphe porte le numéro 24, écrit dans la ligne ; c'est le
premier numéro de paragraphe qui dépasse la dernière page du brouillon
(vue 299, § 23). Le feuillet s'arrête sur « dont une extrémité » ; la
suite n'est pas dans ce lot.}

\end{document}
